\def\linkdethi{https://toanvip307.blogspot.com/} %Đường kinh cần dẫn tới
\hienmaQR
%%%%%%%%%%%%%%%%%%%%%%%% Tiêu đề 3 tham số
\begin{dethi}
	{Huyện Bình Chánh}
	{ĐỀ THAM KHẢO}
	{Kỳ thi tuyển sinh vào lớp 10}
\end{dethi}
\setcounter{EX}{0} 
\subsubsection{Bài tập chọn đáp án đúng}
\Opensolutionfile{ans}[ans/ansTN-Dung]
\begin{ex}%[BG12new-4in1, Trần Hoà]%[2D1H1-1]
	Các khoảng đồng biến của hàm số $y=x^4-8x^2-4$ là
	\choice
	{$(-\infty;-2)$ và $(0;2)$}
	{$(-2;0)$ và $(0;+\infty)$}
	{\True $(-2;0)$ và $(2;+\infty)$}
	{$(-\infty;-2)$ và $(2;+\infty)$}
	\loigiai{Ta có tập xác định của hàm số là $\mathscr{D}=\mathbb{R}$ và $y'=4x^3-16x$ nên $y'=0\Leftrightarrow \hoac{&x=0\\&x=-2\\&x=2.}$\\
		Khi đó $y'\ge 0\Leftrightarrow 4x^3-16x\ge 0\Leftrightarrow \hoac{&x\ge 0\\&-2\le x\le 0.}$\\
		Vậy hàm số đồng biến trên các khoảng $(-2;0)$ và $(2;+\infty)$.
	}
\end{ex}
\begin{ex}%[BG12new-4in1, Trần Hoà]%[2D1H1-1]
	Cho hàm số $y=f(x)$ có đạo hàm $f'(x)=x^2(x^2-1),\,\forall x\in \mathbb{R}$. Hàm số $y=2f(-x)$ đồng biến trên khoảng
	\choice
	{$\left(2;+\infty\right)$}
	{$\left(-\infty;-1\right)$}
	{$(0;2)$}
	{\True $\left(-1;1\right)$}
	\loigiai{
		Hàm số $y=g(x)=2f(-x)$ có đạo hàm $y'=-2f'(-x)=-2x^2(x^2-1)$. Suy ra
		$$y'=0\Leftrightarrow -2x^2(x^2-1)=0\Leftrightarrow \hoac{&x=0\\&x=\pm 1.}$$
		Bảng biến thiên của hàm số $y=g(x)=2f(-x)$ là
		\begin{center}
			\begin{tikzpicture}[scale=1, font=\footnotesize, line join=round, line cap=round, >=stealth]
				\tkzTabInit[nocadre=false,lgt=1.5,espcl=1.5,deltacl=0.6]
				{$x$/.7,$g'(x)$/.7,$g(x)$/1.5}
				{$-\infty$,$-1$,$0$,$1$,$+\infty$}
				\tkzTabLine{,-,z,+,z,+,z,-,} 
				\tkzTabVar{+/,-/,R,+/,-/}
			\end{tikzpicture}
		\end{center}
		Vậy suy ra hàm số $y=g(x)=2f(-x)$ đồng biến trên khoảng $\left(-1;1\right)$.
	}
\end{ex}
\Closesolutionfile{ans}





\subsubsection{Bài tập đúng sai}
\Opensolutionfile{ans}[ans/ansTN-DungSai]
\begin{ex}%[BG12new-4in1, Trần Hoà]%[2D1V1-1]
	Cho hàm số $y=\ln (1+x)$. Xét tính đúng sai của mỗi mệnh đề sau
	\choiceTF[t]
	{Hàm số có tập xác định là $\mathbb{R}$}
	{\True Hàm số đồng biến trên tập xác định}
	{\True Tiếp tuyến của đồ thị vuông góc với đường phân giác góc phần tư thứ hai cắt hai trục toạ độ tại hai điểm phân biệt}
	{$\dfrac{1}{y'(0)}+\dfrac{1}{y'(1)}+\dfrac{1}{y'(2)}+\ldots+\dfrac{1}{y'(2024)}=2025\cdot 2026$}
	\loigiai{
		
		\begin{itemchoice}
			\itemch Sai. Tập xác định $\mathscr{D}=(-1;+\infty)$.
			\itemch Đúng. Ta có $y'=\dfrac{(1+x)'}{(1+x)}=\dfrac{1}{1+x}>0,\,\forall x\in (-1;+\infty)$ nên hàm số đồng biến trên tập xác định.
			\itemch Đúng. Đường phân giác góc phần tư thứ hai có phương trình $y=-x$.\\
			Gọi $k$ là hệ số góc của tiếp tuyến cần tìm, ta có $k=1$.\\
			Hoành độ tiếp điểm là nghiệm phương trình $y'=0\Leftrightarrow \dfrac{1}{1+x}=1\Rightarrow x=0\Rightarrow A(0;1)$.\\
			Phương trình tiếp tuyến là $y=1(x-0)+1\Leftrightarrow y=x+1$.\\
			Đường thẳng cắt hai trục toạ độ tại hai điểm phân biệt.
			\itemch Sai. Ta có
			{\allowdisplaybreaks
				\begin{eqnarray*}
					&&\dfrac{1}{y'(0)}+\dfrac{1}{y'(1)}+\dfrac{1}{y'(2)}+\ldots+\dfrac{1}{y'(2024)}\\
					&=&
					1+2+3+\ldots +2025=\dfrac{2025\cdot 2026}{2}.
			\end{eqnarray*}}	
			
		\end{itemchoice}
	}	
\end{ex}
\begin{ex}%[BG12new-4in1, Trần Hoà]%[2D1V1-1]
	Cho hàm số $y=x+\sqrt{4-x^2}$. Xét tính đúng sai của mỗi mệnh đề sau
	\choiceTF[t]
	{\True Hàm số có tập xác định là $\mathscr{D}=[-2;2]$}
	{$y'=1+\dfrac{x}{\sqrt{4-x^2}}$}
	{Khoảng đồng biến trên khoảng $(-2;2)$}
	{Phương trình $x+\sqrt{4-x^2}=3$ có hai nghiệm phân biệt}
	\loigiai{
		
		\begin{itemchoice}
			\itemch Đúng. Điều kiện: $4-x^2\ge 0\Leftrightarrow -2\le x\le 2 $. Tập xác định $\mathscr{D}=[-2;2]$.
			\itemch Sai. Ta có $y'=1+\dfrac{(4-x^2)'}{2\sqrt{4-x^2}}=1-\dfrac{x}{\sqrt{4-x^2}}$.
			\itemch Sai. Ta có $y'=0\Leftrightarrow \dfrac{x}{\sqrt{4-x^2}}=1\Leftrightarrow x=\sqrt{4-x^2}\Leftrightarrow\heva{&x\ge 0\\&x=\pm \sqrt{2}}\Leftrightarrow x=\sqrt{2}.$\\
			Bảng xét dấu
			\begin{center}
				\begin{tikzpicture}
					\tkzTabInit[nocadre=false,lgt=1,espcl=2.5,deltacl=0.6]%[nocadre=true\fasle(không\có) kẻ viền bảng ,
					%lgt=1,espcl=2]:Độ rộng cột thứ nhât, các cột tiếp theo
					{$x$ /.7,$y'$ /.7,$y$ /2}{$-2$,$\sqrt{2}$,$2$}
					%số dòng của x(1), y'(1), y(2), các giá trị điền ở dòng
					\tkzTabLine{,+,$0$,-,}%điền dấu, giá trị y'
					\tkzTabVar{-/ $0$ ,+/$2\sqrt{2}$,-/$0$} %-/ là giá trị ở dưới thấp
				\end{tikzpicture}
			\end{center}
			Hàm số đồng biến trên khoảng $(-2;\sqrt{2})$ và nghịch biến trên khoảng $(\sqrt{2};2)$. 
			\itemch Sai. Từ bảng biến thiên ta thấy đường thẳng $y=3$ không cắt đồ thị nên phương trình vô nghiệm. 		
		\end{itemchoice}
	}	
\end{ex}
\Closesolutionfile{ans}






\subsubsection{Bài tập điền khuyết}
\Opensolutionfile{ans}[ans/ansTN-DienKhuyet]
\begin{ex}%[BG12new-4in1, Trần Hoà]%[2D1H1-1]
	Hàm số $y=x^4+4x^2+3$. Có bao nhiêu giá trị $x$ nguyên thuộc khoảng đồng biến của hàm số và $x\in [-10;10]$?
	\shortans[]{$-9$}
	\loigiai{
		Ta có $y'=4x^3+8x=4x(x^2+2)$; $y'=0 \Leftrightarrow x=0$.\\
		Bảng biến thiên:
		\begin{center}
			\begin{tikzpicture}
				\tkzTabInit[nocadre=false,lgt=1,espcl=2.5]
				{$x$ /0.6,$y'$ /0.6,$y$ /1.7}
				{$-\infty$,$0$,$+\infty$}
				\tkzTabLine{,-,$0$,+,}
				\tkzTabVar{+/$+\infty$, -/$3$,+/$+\infty$}
			\end{tikzpicture}
		\end{center}
		Suy ra hàm số đồng biến trên khoảng $(0;+\infty )$.\\
		Kết hợp với điều kiện $x$ nguyên thuộc đoạn $[-10;10]$, ta có $x\in\{1;2;\ldots;10\}$.
	}
\end{ex}
\begin{ex}%[BG12new-4in1, Trần Hoà]%[2D1H1-1]
	Hàm số $y=x^4+4x^2+3$. Có bao nhiêu giá trị $x$ nguyên thuộc khoảng đồng biến của hàm số và $x\in [-10;10]$?
	\shortans[]{$-9$}
	\loigiai{
		Ta có $y'=4x^3+8x=4x(x^2+2)$; $y'=0 \Leftrightarrow x=0$.\\
		Bảng biến thiên:
		\begin{center}
			\begin{tikzpicture}
				\tkzTabInit[nocadre=false,lgt=1,espcl=2.5]
				{$x$ /0.6,$y'$ /0.6,$y$ /1.7}
				{$-\infty$,$0$,$+\infty$}
				\tkzTabLine{,-,$0$,+,}
				\tkzTabVar{+/$+\infty$, -/$3$,+/$+\infty$}
			\end{tikzpicture}
		\end{center}
		Suy ra hàm số đồng biến trên khoảng $(0;+\infty )$.\\
		Kết hợp với điều kiện $x$ nguyên thuộc đoạn $[-10;10]$, ta có $x\in\{1;2;\ldots;10\}$.
	}
\end{ex}
\Closesolutionfile{ans}


\newpage
\indapanT{8}{ans/ansTN-Dung}{BÀI TẬP CHỌN ĐÁP ÁN}
\indapanTF{2}{ans/ansTN-DungSai}{BÀI TẬP ĐÚNG SAI}
\indapanSA{4}{ans/ansTN-DienKhuyet}{BÀI TẬP ĐIỀN KHUYẾT}
%%%%%%%%%%%%%%
\label{de\thedeso} %nhãn đề
